% test-001.tex -- tagged-PDF test de los cuatro grandes
%
% Compilar:  lualatex-dev test-001.tex
% Validar:   show-pdf-tags --xml test-001.pdf | rnv-wrapp
%            verapdf --flavour ua2 --format text test-001.pdf
%
% Test de prueba para \spnum, \spunit, \spnewunit, \spunitalias, \spmoney
% y \spshort: las formas de cada comando y sus llaves. Los "doc:" son
% lecturas reales de NVDA 2026.3 con MathCAT 0.7.8-alpha.1. En las unidades, el
% plural (metros, segundos) es lectura propia de MathCAT, no del paquete.
\DocumentMetadata
  {
    lang=es-CL,
    pdfversion=2.0,
    pdfstandard=ua-2,
    tagging=on,
    tagging-setup={math/setup={mathml-SE}},
  }
\documentclass{article}
\usepackage[spanish]{babel}
\usepackage{unicode-math}
\usepackage{spanish-mathml-intent}
\usepackage{hyperref}
\hypersetup
  {
    pdfdisplaydoctitle = true,
    pdftitle           = {test-001: spnum, spunit, spmoney y spshort},
  }
\begin{document}

\section*{A. El comando \texttt{\textbackslash spnum}}

% doc: «más 23»
1. Entero positivo: $\spnum{+23}$.

% doc: «menos 23»
2. Entero negativo: $\spnum{-23}$.

% doc: «3 coma 14»
3. Decimal con coma: $\spnum{3,14}$.

% doc: «3 punto 14»
4. Decimal con punto: $\spnum{3.14}$.

% doc: «3 coma 1 con 4 periódico»
5. Decimal periódico: $\spnum{3,1;4}$.

% doc: «3 coma 4 periódico»
6. Periódico puro: $\spnum{3;4}$.

% doc: «3 coma 14 por 10 al cubo»
7. Notación científica: $\spnum{3,14e3}$.

% doc: «100000»
8. Separación de cifras: $\spnum{100000}$.

% doc: «100000»
9. Sin separación automática: $\spnum[cifras-sep=false]{10\,00\,00}$.

% doc: «23 negativo»
10. Lectura del signo: $\spnum[read-sign=clear]{-23}$.

% doc: «3 punto 4 periódico»
11. Periódico puro con punto: $\spnum[read-period-sep=punto]{3;4}$.

% doc: «3 coma 14 por 10 al cubo»
12. Símbolo de la notación: $\spnum[notation-sym=times]{3,14e3}$.

\section*{B. El comando \texttt{\textbackslash spunit}}

% doc: «23 metros»
13. Unidad con número: $\spunit{23 m}$.

% doc: «metro»
14. Unidad sola: $\spunit{m}$.

% doc: «23 metros segundos al cuadrado»
15. Producto de unidades: $\spunit{23 m*s^{2}}$.

% doc: «23 metros por segundo al cuadrado»
16. Unidades en fracción lineal: $\spunit{23 m/s^{2}}$.

% doc: «90 kilómetros por hora»
17. Velocidad: $\spunit{90 km/h}$.

% doc: «25 grados Celsius»
18. Temperatura: $\spunit{25 °C}$.

% doc: «100 kilogramos»
19. Masa: $\spunit{100 kg}$.

% doc: «23 metros segundos al cuadrado»
20. Producto con punto: $\spunit[print-cdot=true]{23 m*s^{2}}$.

% doc: «23 metros por segundos al cuadrado»
21. Lectura del producto: $\spunit[read-cdot=true]{23 m*s^{2}}$.

% doc: «23 metros partido por segundos al cuadrado»
22. Lectura de la barra: $\spunit[read-slash=partido-por]{23 m/s^{2}}$.

% doc: «3 coma 14 por 10 al cubo metros»
23. Llaves de \texttt{\textbackslash spnum}: $\spunit[notation-sym=times]{3,14e3 m}$.

\section*{C. Los comandos \texttt{\textbackslash spnewunit} y \texttt{\textbackslash spunit}alias}

\spnewunit{pers}{personas}
\spunitalias{caudal}{m^{3}/s}

% doc: «100 personas»
24. Nueva unidad: $\spunit{100 pers}$.

% doc: «50 metros al cubo por segundo»
25. Nuevo alias: $\spunit{50 caudal}$.

\section*{D. El comando \texttt{\textbackslash spmoney}}

% doc: «500 pesos»
26. Moneda por defecto: $\spmoney{500}$.

% doc: «menos 500 pesos»
27. Cifra negativa: $\spmoney{-500}$.

% doc: «1000000 de pesos»
28. Con separación de cifras: $\spmoney{1000000}$.

% doc: «500 dólares»
29. Alias de moneda: $\spmoney{500 dolares}$.

% doc: «500 dólares»
30. Alias con tilde: $\spmoney{500 dólares}$.

% doc: «500 euros»
31. Euros: $\spmoney{500 euro}$.

% doc: «500 libras»
32. Libras: $\spmoney{500 libra}$.

% doc: «500 pesos chilenos»
33. Código ISO: $\spmoney{500 CLP}$.

% doc: «500 dólares estadounidenses»
34. Código ISO en dólares: $\spmoney{500 USD}$.

% doc: «500 euros»
35. Moneda por defecto distinta: $\spmoney[currency=euro]{500}$.

% doc: «100 dólares con 50 centavos»
36. Subunidades: $\spmoney[sub-units=true]{100,5 dolares}$.

% doc: «500 pesos»
37. Símbolo del modo texto: $\spmoney[money-math=false]{500}$.

\section*{E. El comando \texttt{\textbackslash spshort}}

% doc: «doctora»
38. Abreviatura con letra volada: \spshort{dr.a}.

% doc: «doctoras»
39. Abreviatura plural: \spshort{dr.as}.

% doc: «profesora»
40. Abreviatura de título: \spshort{prof.a}.

% doc: «número»
41. Abreviatura de número: \spshort{n.o}.

% doc: «organización de las naciones unidas»
42. Acrónimo: \spshort{onu}.

% doc: «ministerio de educación»
43. Sigla: \spshort{mineduc}.

% doc: «primera»
44. Ordinal femenino: \spshort{1.a}.

% doc: «primero»
45. Ordinal masculino: \spshort{1.o}.

% doc: «doctora»
46. Con versalitas: \spshort[small-caps=true]{dr.a}.

% doc: «octavo»
47. Abreviatura nueva: \spshort[read-arg={octavo}]{8.\textsuperscript{\emph{avo}}}.

\section*{F. Casos adicionales de \texttt{\textbackslash spunit} y \texttt{\textbackslash spnum}}

% doc: «9 coma 8 kilogramos metros por segundo al cuadrado»
48. Producto y cociente con exponente: $\spunit{9,8 kg*m/s^{2}}$.

% doc: «12 kilómetros al cuadrado»
49. Exponente en una sola unidad: $\spunit{12 km^{2}}$.

% doc: «1 coma 5 por 10 elevado a menos 3»
50. Exponente negativo en notación científica: $\spnum{1,5e-3}$.

\section*{G. Unidades de \texttt{\textbackslash spunit} de la versión 1.5}

% doc: «10 micrómetros»
51. Prefijo micro pegado a la unidad: $\spunit{10 µm}$.

% doc: «un microgramo»
52. Prefijo micro con una unidad: $\spunit{1 µg}$.

% doc: «25 grados Celsius»
53. Símbolo compacto de grados: $\spunit{25 ℃}$.

% doc: «25 metros cuadrados»
54. Símbolo compacto de área: $\spunit{25 ㎡}$.

% doc: «10 metros por segundo»
55. Símbolo compacto con barra: $\spunit{10 ㎧}$.

% doc: «5 kilovatios hora»
56. Unidad de energía: $\spunit{5 kWh}$.

% doc: «760 milímetros de mercurio»
57. Unidad de presión: $\spunit{760 mmHg}$.

% doc: «2 henrios»
58. Henrio: $\spunit{2 H}$.

% doc: «3 néperes»
59. Néper: $\spunit{3 Np}$.

\end{document}
